\problemstart
\addcontentsline{toc}{section}{2.1　斜面上の束縛運動と摩擦}

\begin{problemBox}{問題 2.1\quad 斜面上の束縛運動と摩擦\hspace{1.2em}{\small\mdseries 〔基本〕}}
水平面に対して角度 $0 < \theta < \frac{\pi}{2}$ だけ傾いた粗い斜面上に，質量 $m$ の小物体を置く。
\par
斜面に沿って下る向きに $x$ 軸をとり，斜面に垂直で物体から離れる向きに $y$ 軸をとる。重力加速度の大きさを $g$，静止摩擦係数を $\mu$，動摩擦係数を $\mu'<\mu$ とする。
\par
時刻 $t = 0$ において，原点 $x = 0$ に静止させた物体から静かに手を離す。以下の問いに答えよ。

\begin{enumerate}
  \item 物体が自発的に滑り始めるための傾斜角 $\theta$ に関する条件を，静止摩擦係数 $\mu$ を用いて表せ。

  \item 角度 $\theta$ が (1) で求めた条件を満たしているとき，物体は滑り始める。滑り始めた後の $x$ 方向の運動方程式を，速度 $v(t) = \frac{dx}{dt}$ に関する1階微分方程式として表せ。また，この運動における物体の加速度を求めよ。

  \item 初期条件 $v(0) = 0$ および $x(0) = 0$ のもとで，(2) の微分方程式を解き，時刻 $t$ における物体の速度 $v(t)$ および位置 $x(t)$ を求めよ。
\end{enumerate}
\end{problemBox}

\solutionheading
\begin{solutionparts}

\solutionpart{(1)}
物体が斜面上で静止しているとき，物体に作用する力は，重力 $mg$，斜面からの垂直抗力 $N$，および静止摩擦力 $f$ です。 $y$ 方向の力のつり合いより，垂直抗力 $N$ は以下のように求まります。
\begin{equation*}
N - m g \cos \theta = 0
\end{equation*}
\begin{equation*}
N = m g \cos \theta
\end{equation*}
一方，$x$ 方向の力のつり合いを考えると，重力の斜面平行成分 $m g \sin \theta$ と静止摩擦力 $f$ がつり合っています。
\begin{equation*}
f = m g \sin \theta
\end{equation*}
物体が静止し続けるためには，この静止摩擦力 $f$ が最大静止摩擦力 $f_{\text{max}} = \mu N$ 以下である必要があります。
\begin{equation*}
f \le \mu N
\end{equation*}
したがって，物体が自発的に滑り出すための条件は，重力の斜面平行成分が最大静止摩擦力を超えること，すなわち $f > \mu N$ となることです。
\begin{equation*}
m g \sin \theta > \mu m g \cos \theta
\end{equation*}
両辺を $m g \cos \theta > 0$ で割ることで，求める条件が得られます。
\begin{equation*}
\tan \theta > \mu
\end{equation*}
物理的には，斜面の傾き $\theta$ が摩擦角 $\arctan \mu$ よりも大きくなったときに，物体は滑り始めることを意味しています。

\solutionpart{(2)}
物体が滑り始めた後は，物体に作用する摩擦力は動摩擦力 $f'$ となります。動摩擦力の大きさは，垂直抗力 $N$ を用いて以下のように表されます。
\begin{equation*}
f' = \mu' N = \mu' m g \cos \theta
\end{equation*}
動摩擦力は物体の運動を妨げる向きに働くため，$x$ 軸の負の向きに作用します。 したがって，滑り出した後の $x$ 方向の運動方程式は次のようになります。
\begin{equation*}
m \frac{dv}{dt} = m g \sin \theta - \mu' m g \cos \theta
\end{equation*}
両辺を質量 $m$ で割ることで，速度 $v(t)$ に関する1階微分方程式が得られます。
\begin{equation*}
\frac{dv}{dt} = g (\sin \theta - \mu' \cos \theta)
\end{equation*}
この式の右辺は時間 $t$ に依存しない定数であるため，物体の加速度 $a$ は一定となります。したがって，求める加速度 $a$ は以下の通りです。
\begin{equation*}
a = g (\sin \theta - \mu' \cos \theta)
\end{equation*}
(1) の条件 $\tan \theta > \mu$ かつ $\mu > \mu'$ より，$\sin \theta - \mu' \cos \theta > 0$ が成り立つため，加速度は正の値となり，物体は斜面を下向きに加速していきます。

\solutionpart{(3)}
(2) で得られた微分方程式は，両辺を時間 $t$ で直接積分することで解くことができます。
\begin{equation*}
\frac{dv}{dt} = a
\end{equation*}
両辺を $t$ について積分すると，積分定数を $C_1$ として，以下のようになります。
\begin{equation*}
v(t) = a t + C_1
\end{equation*}
初期条件 $v(0) = 0$ を代入すると， $C_1 = 0$ と定まります。 したがって，時刻 $t$ における物体の速度 $v(t)$ は以下のように求まります。
\begin{equation*}
v(t) = g (\sin \theta - \mu' \cos \theta) t
\end{equation*}
次に，位置 $x(t)$ は速度の定義 $v(t) = \frac{dx}{dt}$ より，速度を時間 $t$ でさらに積分することで得られます。
\begin{equation*}
x(t) = \int v(t) dt = \frac{1}{2} g (\sin \theta - \mu' \cos \theta) t^2 + C_2
\end{equation*}
ここで $C_2$ は積分定数です。初期条件 $x(0) = 0$ を代入すると， $C_2 = 0$ と定まります。 したがって，時刻 $t$ における物体の位置 $x(t)$ は以下のようになります。
\begin{equation*}
x(t) = \frac{1}{2} g (\sin \theta - \mu' \cos \theta) t^2
\end{equation*}
したがって，物体は加速度 $g (\sin \theta - \mu' \cos \theta)$ の等加速度運動を行います。

\end{solutionparts}
